% !TEX root = ../../tree_notes.tex
% =====================================================================
%  3.3  Breadth-First / Level Order                         sec:c03-bfs
%  Exports: def:c03-level  def:c03-frontier  thm:c03-bfs-levels
%           thm:c03-bfs-shortest  lst:c03-bfs  ex:c03-bfs-run
%           prob:c03-nary-level  prob:c03-zigzag  prob:c03-largest-row
% =====================================================================
\section{Breadth-First / Level Order}\label{sec:c03-bfs}

\begin{previouslybox}
\begin{enumerate}[nosep, leftmargin=1.6em]
  \item What does the explicit stack in inorder traversal remember, and what is the
        breadth-first analogue we build today? (\defref{c03-explicit-stack})
  \item Morris reached what space bound, and why can level order not? (\thmref{c03-morris-linear})
  \item On $R$, list the depth-first orders once more -- which of them already looks
        ``row by row''? (\secref{c03-dfs}, \figref{c03-running})
\end{enumerate}
\end{previouslybox}

\begin{topicintro}
Depth-first plunges to a leaf before backing off; breadth-first does the opposite -- it
finishes every node at depth $d$ before touching depth $d{+}1$. The resulting sequence is
\emph{row by row}, and it is produced by a queue instead of a stack. This single change of
data structure turns the same walk into the tool for everything shaped like ``by level'':
per-row statistics, the right-side view, and (on general graphs) shortest hop-counts. The
one idea to internalise is the \term{level frontier}: snapshot \code{q.size()} once to know
exactly how many nodes belong to the current row.
\textbf{Prerequisites.} \secref{c03-dfs} (the DFS family, for contrast); FIFO queue basics.

\textbf{After this section you can}
\begin{itemize}[nosep]
  \item emit a tree level-by-level with the queue + frontier-size template;         \hfill\lv{2}
  \item adapt it to zigzag, bottom-up, and per-row aggregate variants;              \hfill\lv{3}
  \item handle N-ary nodes and \code{next}-pointer rewiring inside the same loop;   \hfill\lv{3}--\lv{4}
\end{itemize}
\end{topicintro}

\begin{recognitionbox}[Pattern card -- breadth-first / level order]
\reckey{``by level / row / from top to bottom'', ``zigzag level order'', ``right side
view'', ``the largest value in each row'', ``average per level'' -- the answer is grouped
by \emph{depth}.}
\recuse{
  \item the output must be organised \emph{per depth} (a list of lists, one row each);
  \item a per-row statistic (max, mean, the visible right edge) is requested;
  \item hop-count / shortest-level questions, where first-visit order matters.
}
\recnotuse{
  \item the answer needs a fixed \emph{subtree} order (pre/in/post) -- \secref{c03-dfs};
  \item memory is capped at $\bigO{1}$ -- level order needs a whole frontier
        ($\bigO{w}$); use Morris (\secref{c03-iterative}) if a DFS order will do;
  \item a plain ``does a path with property $P$ exist'' -- a Chapter~2 recursion may beat BFS.
}
One-line test: \emph{would the answer change if I regrouped nodes by depth? If the problem
speaks in rows, you are here.}
\end{recognitionbox}

\subsection{Intuition}

Picture the tree cut into horizontal bands. Breadth-first sweeps one band at a time, and
the queue always holds exactly ``the next band, waiting.'' The subtle move -- the one every
correct level-order solution makes -- is to \emph{freeze} the queue length at the top of
each outer iteration: that number is how many nodes are in the current row, and processing
precisely that many pops the whole row while enqueuing the next.

\begin{figure}[H]
\centering
\begin{tikzpicture}[diagram, scale=0.82, transform shape]
  \node[troot] (r) at (0,0) {$1$};
  \node[tnode] (a) at (-1.7,-1.4) {$2$};
  \node[tnode] (b) at (1.7,-1.4) {$3$};
  \node[tleaf] (c) at (-2.6,-2.8) {$4$};
  \node[tleaf] (d) at (-0.8,-2.8) {$5$};
  \node[tnode] (e) at (1.7,-2.8) {$6$};
  \draw[tedge] (r) -- (a); \draw[tedge] (r) -- (b);
  \draw[tedge] (a) -- (c); \draw[tedge] (a) -- (d);
  \draw[tedge] (b) -- (e);
  \draw[<-, accentorange, line width=1pt] (-3.6,-0.05) -- (3.4,-0.05) node[right, font=\scriptsize, text=accentorange] {level 0};
  \draw[<-, primaryblue,  line width=1pt] (-3.6,-1.45) -- (3.4,-1.45) node[right, font=\scriptsize, text=primaryblue] {level 1};
  \draw[<-, successgreen,line width=1pt] (-3.6,-2.85) -- (3.4,-2.85) node[right, font=\scriptsize, text=successgreen] {level 2};
  \node[annot, anchor=north west, align=left, font=\scriptsize] at (4.2,0.2)
    {start of row 1:\\ queue $=[\,2\;\;3\,]$,\\ snapshot \code{sz}$=2$\\ $\Rightarrow$ pop 2 nodes};
\end{tikzpicture}
\caption{Level order on $R$ sweeps one band at a time. At the start of processing level
$1$ the queue holds exactly that row $\{2,3\}$; freezing \code{sz}$=2$ tells the inner loop
to pop two nodes (emitting $2,3$) and enqueue their children $4,5,6$ as the next row.
Output: $1$ / $2,3$ / $4,5,6$.}\label{fig:c03-bfs}
\end{figure}

\subsection{Definitions and Notation}

\begin{definition}{Level (breadth-first) order}{c03-level}
The \term{level order} of a rooted tree lists nodes by increasing depth, breaking ties
left-to-right within a level: all of level $0$, then level $1$, and so on. It is produced
by \term{breadth-first search} (BFS), which explores the frontier nearest the root first.
\end{definition}

\begin{definition}{Level frontier}{c03-frontier}
During BFS, at the instant an outer iteration begins, the queue content is exactly one
whole level; its size is the \term{frontier width} $w$ of that level. The maximum such $w$
over the run is the \emph{width} of the tree and bounds BFS memory.
\end{definition}

\noindent\textbf{In plain words.} A stack reaches \emph{down} (depth-first); a queue reaches
\emph{across} (breadth-first). ``Snapshot the size'' is how you tell a queue, which has no
notion of ``level,'' where one row ends and the next begins.

\subsection{Properties}

\begin{theorem}{One Iteration, One Level}{c03-bfs-levels}
The template \coderef{c03-bfs}, with \code{sz = q.size()} taken before the inner loop,
emits exactly the nodes of one depth per outer iteration, in left-to-right order.
\end{theorem}
\begin{proofbox}
Induction on the outer iteration, i.e.\ on depth $d$. Base $d=0$: the queue starts with
only the root, \code{sz}$=1$, and the inner loop emits it and enqueues precisely its
children (level $1$). Step: assume at the top of iteration $d$ the queue holds exactly
level $d$ (left-to-right) and nothing deeper. \code{sz} freezes that count; the inner loop
pops all $d$-level nodes \emph{before} any newly pushed node (FIFO), emitting them in order
and appending their children -- which are exactly level $d{+}1$ -- to the back, preserving
left-to-right order. So at the top of iteration $d{+}1$ the queue holds level $d{+}1$
only. The invariant holds, giving one level per iteration.
\end{proofbox}

\begin{theorem}{BFS Reaches Shallowest First}{c03-bfs-shortest}
BFS visits nodes in non-decreasing depth; the first time a node is dequeued is via a
minimum-edge path from the root.
\end{theorem}
\begin{proofbox}
By \thmref{c03-bfs-levels} all depth-$d$ nodes are emitted before any depth-$(d{+}1)$
node, so visit order is non-decreasing in depth. A node at depth $d$ is enqueued by a
depth-$(d{-}1)$ parent, hence reached by a $d$-edge path; no shallower path exists (else
its parent would be at depth $<d{-}1$ and it would have been enqueued an iteration
earlier). Thus first-dequeue depth equals true minimum distance.
\end{proofbox}
\begin{insightbox}
On a \emph{tree} every root-to-node path is unique, so \thmref{c03-bfs-shortest} is only
``deeper later''; its real power appears on general graphs (\chref{c07}), where it is the
shortest-path algorithm. We keep it here to justify ``first-touch is optimal'' reasoning
used by the right-side-view and depth variants.
\end{insightbox}

\subsection{The Template}

\begin{lstlisting}[caption={Binary-tree level order -- the frontier-size snapshot}, label={lst:c03-bfs}]
#include <bits/stdc++.h>
using namespace std;

struct TreeNode {
    int val;
    TreeNode *left, *right;
    TreeNode(int x = 0) : val(x), left(nullptr), right(nullptr) {}
};

vector<vector<int>> levelOrder(TreeNode* root) {
    vector<vector<int>> res;
    if (!root) return res;
    queue<TreeNode*> q;
    q.push(root);
    while (!q.empty()) {
        int sz = q.size();                 // FREEZE the frontier: this row's width
        vector<int> level;
        while (sz--) {                     // process exactly `sz` nodes == one level
            TreeNode* n = q.front(); q.pop();
            level.push_back(n->val);
            if (n->left)  q.push(n->left); // next level accumulates at the back
            if (n->right) q.push(n->right);
        }
        res.push_back(level);
    }
    return res;
}
\end{lstlisting}
\begin{complexitybox}
Time $\bigO{n}$: each node is enqueued and dequeued once. Space $\bigO{w}$ where $w$ is the
maximum level width ($\bigO{n}$ worst case -- a perfect tree's last row has $\sim n/2$
nodes). This frontier memory is the price of breadth-first order and cannot be driven to
$\bigO{1}$ the way Morris drives a DFS order.
\end{complexitybox}

\subsection{Worked Example}

\begin{example}{Level order of $R$}{c03-bfs-run}
Freeze-sizes go $1,2,3$: iteration~1 pops $1$ (push $2,3$); iteration~2 pops $2,3$
(push $4,5$ then nothing from $3$'s missing left, then $6$); iteration~3 pops $4,5,6$.
Result \code{\{[1],[2,3],[4,5,6]\}}.
\end{example}
\begin{tracebox}[Queue content at each outer top on $R$]
\begin{itemize}[nosep, leftmargin=1.4em]
  \item top 1: $[\,1\,]$, \code{sz}=1 $\to$ emit $1$; enqueue $2,3$
  \item top 2: $[\,2,3\,]$, \code{sz}=2 $\to$ emit $2,3$; enqueue $4,5,6$ (note $3$'s
        null left enqueues nothing)
  \item top 3: $[\,4,5,6\,]$, \code{sz}=3 $\to$ emit $4,5,6$; enqueue none
  \item queue empty $\to$ stop. Levels: $1$ / $2,3$ / $4,5,6$.
\end{itemize}
\end{tracebox}

\subsection{Variations and Twists}

\begin{itemize}[nosep]
  \item \textbf{Zigzag}: reverse every other emitted row (or fill a preallocated
        \code{level} from both ends).                       \hfill\lv{3}
  \item \textbf{Bottom-up}: run level order, then \code{reverse(res)} -- LC 107.  \hfill\lv{2}
  \item \textbf{Right side view}: push only \code{level.back()} of each row -- LC 199,
        referenced in \secref{c03-rebuild}'s practice notes.  \hfill\lv{2}
  \item \textbf{Per-row statistic}: fold max/sum/avg inside the inner loop --
        \probref{c03-largest-row}, LC 637.                     \hfill\lv{2}
  \item \textbf{N-ary nodes}: replace the two fixed \code{push} calls with a loop over
        \code{children} -- \probref{c03-nary-level}.            \hfill\lv{2}
\end{itemize}

\subsection{Common Pitfalls}\label{sec:c03-bfs-pitfalls}

\begin{warningbox}[Trap 1 -- reading \code{q.size()} inside the inner loop]
The frontier \emph{grows} as you enqueue children, so \code{while (i < q.size())} never
terminates the row correctly (it chases its own tail). Snapshot once \emph{before} the
inner loop.
\end{warningbox}
\begin{lstlisting}[caption={\textcolor{warningred}{WRONG}: moving goalpost -- row never ends}, language=C++, basicstyle=\ttfamily\small]
while (!q.empty()) {
    for (int i = 0; i < (int)q.size(); ++i) {   // q.size() re-read each test!
        auto* n = q.front(); q.pop(); /* ... push children ... */
    }
}
\end{lstlisting}
\begin{lstlisting}[caption={\textcolor{successgreen}{RIGHT}: freeze the width once}, language=C++, basicstyle=\ttfamily\small]
while (!q.empty()) {
    int sz = q.size();                           // one snapshot = this row's width
    while (sz--) { auto* n = q.front(); q.pop(); /* ... push children ... */ }
}
\end{lstlisting}

\begin{warningbox}[Trap 2 -- enqueuing null children]
Pushing \code{n->left} unconditionally puts $\nil$ in the queue, so you dereference it
next iteration (RE) or emit phantom values. Guard both pushes with a non-null test -- and
do \emph{not} rely on a \code{if(!n) continue} patch to hide it.
\end{warningbox}
\begin{lstlisting}[caption={\textcolor{warningred}{WRONG}: nulls pollute the frontier}, language=C++, basicstyle=\ttfamily\small]
q.push(n->left); q.push(n->right);   // may enqueue nullptr
\end{lstlisting}
\begin{lstlisting}[caption={\textcolor{successgreen}{RIGHT}: only real nodes enter the queue}, language=C++, basicstyle=\ttfamily\small]
if (n->left)  q.push(n->left);
if (n->right) q.push(n->right);
\end{lstlisting}

\begin{warningbox}[Trap 3 -- forgetting the empty-tree guard]
\code{q.push(root)} on a null root (LC's ``tree can be empty'' case) makes the first
iteration dereference $\nil$. Return the empty result before touching the queue.
\end{warningbox}
\begin{lstlisting}[caption={\textcolor{warningred}{WRONG}}  , language=C++, basicstyle=\ttfamily\small]
queue<TreeNode*> q; q.push(root);   // root may be nullptr
\end{lstlisting}
\begin{lstlisting}[caption={\textcolor{successgreen}{RIGHT}}, language=C++, basicstyle=\ttfamily\small]
if (!root) return res;               // handle the empty tree up front
queue<TreeNode*> q; q.push(root);
\end{lstlisting}


\subsection{Problem Lab}

One \lv{2} warm-up (N-ary level order -- the template with a children loop) solved
concisely, then two \lv{3} core problems that reshape the emitted rows: zigzag
(reverse alternate rows) and per-row maximum (fold inside the loop).

\solvedproblem{c03-nary-level}{N-ary Tree Level Order Traversal}{LeetCode 429}
  {https://leetcode.com/problems/n-ary-tree-level-order-traversal/}{LC Medium L2}
  {BFS, frontier snapshot, N-ary children loop}

\begin{problemstatement}[N-ary Level Order]
Given an N-ary tree (each node has a \code{vector<Node*> children}), return its level
order traversal -- one list per depth, children in the given left-to-right order.
\end{problemstatement}

\solstep{Understanding the Problem}
Identical to \coderef{c03-bfs} except a node may have any number of children: replace the
two guarded pushes with a single loop over \code{children}. The frontier snapshot still
draws the level boundary.

\solstep{Key Observations}
\begin{observation}
The child-count is irrelevant to BFS -- you enqueue all of them and the frozen \code{sz}
already accounts for the whole next row.
\end{observation}
\begin{observation}
The empty tree is \code{root == nullptr}; guard before pushing (Trap 3).
\end{observation}

\solstep{Solution}
Standard level-order loop; inner body emits \code{n->val} then pushes every child.

\solstep{Implementation}
\begin{lstlisting}[caption={N-ary Level Order (LeetCode 429)}, label={lst:c03-nary-level}]
#include <bits/stdc++.h>
using namespace std;

class Node {
public:
    int val;
    vector<Node*> children;
    Node() {}
    Node(int v) : val(v) {}
};

class Solution {
public:
    vector<vector<int>> levelOrder(Node* root) {
        vector<vector<int>> res;
        if (!root) return res;
        queue<Node*> q; q.push(root);
        while (!q.empty()) {
            int sz = q.size();
            vector<int> level;
            while (sz--) {
                Node* n = q.front(); q.pop();
                level.push_back(n->val);
                for (Node* c : n->children) q.push(c);  // the ONLY change vs binary BFS
            }
            res.push_back(level);
        }
        return res;
    }
};
\end{lstlisting}

\solstep{Trace}
\begin{tracebox}[On LC's sample \code{[1,null,2,3,4,null,null,5,6]}]
Row 0: $[1]$. Row 1: $[2,3,4]$ (children of $1$). Row 2: $[5,6]$ (children of $2$).
Result \code{\{[1],[2,3,4],[5,6]\}}. \ac.
\end{tracebox}

\solstep{Complexity}
\begin{complexitybox}
Time $\bigO{n}$, space $\bigO{w}$ (widest row). Verdict: \ac.
\end{complexitybox}

\begin{insightbox}[Key takeaway]
BFS is \emph{representation-agnostic}: binary or N-ary, only the ``add my children'' step
changes -- the frontier logic is untouched. Compare with \probref{c03-preorder}'s N-ary
preorder (LC 589): same principle on the DFS side.
\end{insightbox}

\solvedproblem{c03-zigzag}{Binary Tree Zigzag Level Order Traversal}{LeetCode 103}
  {https://leetcode.com/problems/binary-tree-zigzag-level-order-traversal/}{LC Medium L3}
  {BFS, conditional row reversal}

\begin{problemstatement}[Zigzag Level Order]
Return level order traversal, but alternating direction: level $0$ left-to-right, level
$1$ right-to-left, level $2$ left-to-right, and so on.
\end{problemstatement}

\begin{attemptgate}[Attempt it yourself first -- \lv{3}, 15 minutes]
A fair attempt: produce plain level order, then fix the direction per row.
\peek{is the \emph{walk} actually different, or only the \emph{listing} of each row?}
\nudge{collect each row normally; keep a boolean that flips every level and reverse the
row when it is \code{false}.}
\fullhint{BFS + \code{bool leftToRight}; \code{if(!leftToRight) reverse(level.begin(),
level.end());} before pushing. The queue order never changes -- only the emitted row. Full
build below.}
\end{attemptgate}

\solstep{Understanding the Problem}
The traversal (what enters the queue and when) is \emph{exactly} \coderef{c03-bfs}; only
the \emph{presentation} of alternate rows is mirrored. Recognising that the BFS machinery
is unchanged is the whole trick -- do not try to traverse right-to-left by reordering
pushes (that breaks the next level).

\solstep{Key Observations}
\begin{observation}
Reversing a collected row is $\bigO{\text{width}}$ and local to one level; it cannot
corrupt the queue because the queue was already filled in canonical order.
\end{observation}
\begin{observation}
Direction is a pure function of level parity -- one flipping boolean suffices, no depth
counter needed.
\end{observation}

\solstep{Think It Through With Me}
\begin{thinkbox}
\deadend{Alternate the \emph{push} order (children right-then-left on odd levels).}{The
queue now holds next-level nodes in mirrored order, so \emph{their} children get mirrored
too, compounding the error every level -- you fix rows $1,3,5$ but break $2,4$. Direction
must live in the output, not the frontier.}
\aha{run canonical BFS untouched; reverse only the vector you are about to store on odd
levels.}
\end{thinkbox}

\solstep{The Approach Ladder}
\begin{center}\renewcommand{\arraystretch}{1.3}
\begin{tabular}{@{}c p{0.235\linewidth} p{0.28\linewidth} l c@{}}
\toprule
\textbf{\#} & \textbf{Approach} & \textbf{Idea} & \textbf{Time} & \textbf{Verdict} \\
\midrule
1 & \lv{2} mirror the pushes & reorder enqueue & $O(n)$ & \textbf{WA} \\
2 & \lv{3} BFS + reverse row & canonical BFS, flip output & $O(n)$ & \ac \\
\bottomrule
\end{tabular}\end{center}

\solstep{Solution}
Plain level-order BFS. Maintain \code{bool ltr=true}; after building \code{level}, if
\code{!ltr} reverse it, push to \code{res}, then flip \code{ltr}.

\solstep{Proof of Correctness}
\begin{proofbox}
By \thmref{c03-bfs-levels} the un-reversed rows are the correct left-to-right level lists.
Reversing row $d$ when $d$ is odd yields right-to-left for those rows and leaves even rows
left-to-right, which is precisely the zigzag specification; the queue state is untouched by
the reversal, so all rows remain correct.
\end{proofbox}

\solstep{Implementation}
\begin{lstlisting}[caption={Zigzag Level Order (LeetCode 103)}, label={lst:c03-zigzag}]
class Solution {
public:
    vector<vector<int>> zigzagLevelOrder(TreeNode* root) {
        vector<vector<int>> res;
        if (!root) return res;
        queue<TreeNode*> q; q.push(root);
        bool ltr = true;                            // level 0 goes left -> right
        while (!q.empty()) {
            int sz = q.size();
            vector<int> level;
            while (sz--) {
                TreeNode* n = q.front(); q.pop();
                level.push_back(n->val);
                if (n->left)  q.push(n->left);
                if (n->right) q.push(n->right);
            }
            if (!ltr) reverse(level.begin(), level.end());   // only the OUTPUT flips
            res.push_back(level);
            ltr = !ltr;
        }
        return res;
    }
};
\end{lstlisting}

\solstep{Trace}
\begin{tracebox}[On $R$]
Level 0: \code{ltr}$=$true $\to [1]$. Level 1 collected $[2,3]$, \code{ltr}$=$false
$\to$ reversed $[3,2]$. Level 2 collected $[4,5,6]$, \code{ltr}$=$true $\to [4,5,6]$.
Result \code{\{[1],[3,2],[4,5,6]\}}. \ac.
\end{tracebox}

\solstep{Complexity}
\begin{complexitybox}
Time $\bigO{n}$ (each node visited once; total reversal work over all rows is
$\bigO{n}$). Space $\bigO{w}$. A deque variant fills odd rows front-first to avoid the
explicit reverse -- same asymptotics, marginally fewer writes.
\end{complexitybox}

\solstep{Verification and Edge Cases}
\begin{center}\renewcommand{\arraystretch}{1.2}
\begin{tabular}{@{}p{0.30\linewidth} p{0.62\linewidth}@{}}
\toprule
\textbf{Case} & \textbf{What the code does} \\
\midrule
Empty tree & returns \{\} at the guard \\
Single node & one row $[\,\mathit{root}\,]$, no reversal \\
Full last level & reversal cost $\bigO{n/2}$ but still linear overall \\
\bottomrule
\end{tabular}\end{center}

\solstep{Pitfalls and Variants}
\begin{warningbox}
Reversing \emph{inside} the inner loop (per node) instead of once per row costs a factor
and confuses indices; collect the whole row, then reverse once. Variant: fill
\code{level[mid $\pm$ j]} from both ends to skip the reverse (a matter of constant
factor, not correctness).
\end{warningbox}

\begin{insightbox}[Key takeaways -- three things to remember]
\begin{enumerate}[nosep, leftmargin=1.6em]
  \item When a problem changes only how a level is \emph{presented}, keep the traversal
        canonical and post-process the row.
  \item ``Don't touch the frontier for a cosmetic requirement'' generalises: the queue/
        stack invariant is fragile; output shaping is free.
  \item Level parity is a $\bigO{1}$-state toggle -- most zigzag bugs come from tracking it
        wrong, not from BFS.
\end{enumerate}
\end{insightbox}

\solvedproblem{c03-largest-row}{Find Largest Value in Each Tree Row}{LeetCode 515}
  {https://leetcode.com/problems/find-largest-value-in-each-tree-row/}{LC Medium L3}
  {BFS with per-row reduction, or DFS-with-depth alternative}

\begin{problemstatement}[Largest Value Per Row]
Given a binary tree, return a vector whose $i$-th element is the maximum node value at
depth $i$. Node values fit in signed 32-bit (may be negative).
\end{problemstatement}

\begin{attemptgate}[Attempt it yourself first -- \lv{3}, 15 minutes]
A fair attempt: sweep levels and take each row's max.
\peek{does a per-row reduction need anything beyond the frontier loop?}
\nudge{fold \code{rowMax = max(rowMax, n->val)} inside the inner loop; reset per row; use
\code{INT\_MIN} as the seed since values can be negative.}
\fullhint{level-order BFS, one \code{curMax} per row reinitialised to \code{INT\_MIN}.
A DFS that indexes \code{res[depth]} is the space-friendly twin. Full build below.}
\end{attemptgate}

\solstep{Understanding the Problem}
This is a per-row \emph{aggregate}, not a per-row \emph{listing} -- the same frontier loop
reduced with \code{max} instead of collected. The negative-value range is the one real
hazard (seeding \code{curMax} at $0$ fails a row of all-negatives).

\solstep{Key Observations}
\begin{observation}
The maximum of a row is computed by folding \code{max} over the inner loop's nodes; no
extra storage per row is needed (contrast \probref{c03-zigzag}).
\end{observation}
\begin{observation}
Values can be negative, so the per-row seed must be \code{INT\_MIN}, never $0$.
\end{observation}
\begin{observation}
A DFS reaches the same answer by keeping \code{res[depth]} and updating it on each visit --
$\bigO{h}$ stack instead of $\bigO{w}$ queue.
\end{observation}

\solstep{Think It Through With Me}
\begin{thinkbox}
\deadend{Recompute each row's max by a fresh DFS to depth $d$ for every $d$.}{Correct but
$\bigO{n \cdot h}$ -- you re-walk upper levels for each row.}
\deadend{Init \code{curMax}$=0$.}{A row like $\{-5,-2\}$ yields max $0$, wrong; the
negative range is a deliberate trap.}
\aha{One BFS sweep, one fold of \code{max} per row, seeded at \code{INT\_MIN}.}
\end{thinkbox}

\solstep{The Approach Ladder}
\begin{center}\renewcommand{\arraystretch}{1.3}
\begin{tabular}{@{}c p{0.235\linewidth} p{0.28\linewidth} l c@{}}
\toprule
\textbf{\#} & \textbf{Approach} & \textbf{Idea} & \textbf{Time} & \textbf{Verdict} \\
\midrule
1 & \lv{2} per-row depth DFS & re-walk for each level & $O(n\cdot h)$ & \tle \\
2 & \lv{3} BFS fold & one sweep, \code{max} per row & $O(n)$, $O(w)$ & \ac \\
3 & \lv{3} DFS + depth index & \code{res[depth]} update & $O(n)$, $O(h)$ & \ac \\
\bottomrule
\end{tabular}\end{center}
Rungs 2 and 3 tie asymptotically (which is ``better'' depends on whether $h$ or $w$
dominates); we build BFS to stay inside this section's tool.

\solstep{Solution}
Level-order BFS; for each row run \code{sz} pops, doing \code{curMax = max(curMax,
n->val)} and enqueueing children; push \code{curMax} after the row.

\solstep{Proof of Correctness}
\begin{proofbox}
By \thmref{c03-bfs-levels} iteration $d$ processes exactly the depth-$d$ nodes, so the
fold $\max$ over those pops equals the row maximum. \code{INT\_MIN} is the identity of
\code{max} over a non-empty set and is $\le$ every legal value, so a negative-only row is
handled correctly; every row is non-empty by construction (BFS stops when there are none).
\end{proofbox}

\solstep{Implementation}
\begin{lstlisting}[caption={Largest Value in Each Row (LeetCode 515)}, label={lst:c03-largest-row}]
class Solution {
public:
    vector<int> largestValues(TreeNode* root) {
        vector<int> res;
        if (!root) return res;
        queue<TreeNode*> q; q.push(root);
        while (!q.empty()) {
            int sz = q.size();
            int curMax = INT_MIN;                 // identity for max; handles negatives
            while (sz--) {
                TreeNode* n = q.front(); q.pop();
                curMax = max(curMax, n->val);
                if (n->left)  q.push(n->left);
                if (n->right) q.push(n->right);
            }
            res.push_back(curMax);
        }
        return res;
    }
};
\end{lstlisting}

\solstep{Trace}
\begin{tracebox}[On $R$]
Row 0 \code{max(INT\_MIN,1)}$=1$; Row 1 over $\{2,3\}$ $=3$; Row 2 over $\{4,5,6\}$ $=6$.
Result \code{[1,3,6]}. \ac.
\end{tracebox}

\solstep{Complexity}
\begin{complexitybox}
Time $\bigO{n}$, space $\bigO{w}$ (queue). The DFS twin is $\bigO{n}$ / $\bigO{h}$, which
wins memory on a wide-but-shallow tree and loses on a tall-thin one.
\end{complexitybox}

\solstep{Verification and Edge Cases}
\begin{center}\renewcommand{\arraystretch}{1.2}
\begin{tabular}{@{}p{0.30\linewidth} p{0.62\linewidth}@{}}
\toprule
\textbf{Case} & \textbf{What the code does} \\
\midrule
All-negative values & \code{INT\_MIN} seed yields the true (negative) row max \\
Single node & returns \code{[\,root->val\,]} \\
Value $=\texttt{INT\_MIN}$ & seed equals value; \code{max} still returns it -- safe \\
\bottomrule
\end{tabular}\end{center}

\solstep{Pitfalls and Variants}
\begin{warningbox}
Seeding \code{curMax}$=0$ (or \code{curMax}$=$ first child value read lazily) is the classic
miss on all-negative rows. Variant LC 199 (right side view) replaces the fold with
``keep \code{level.back()}''; LC 637 (row averages) folds \code{sum} and divides by the
snapshot \code{sz} -- both are this loop with a different reducer.
\end{warningbox}

\begin{insightbox}[Key takeaways -- three things to remember]
\begin{enumerate}[nosep, leftmargin=1.6em]
  \item ``Per level, compute $X$'' is one template: swap the reducer (\code{max},
        \code{sum}, collect) to serve a dozen LC problems.
  \item Reduction identities (\code{INT\_MIN} for max, $0$ for sum) must respect the value
        domain -- the negative case is where these die.
  \item BFS ($\bigO{w}$) and DFS-with-depth ($\bigO{h}$) are interchangeable for per-row
        answers; pick by which of width/height dominates.
\end{enumerate}
\end{insightbox}

\subsection{Practice Problems}

\begin{stuckbox}
\begin{itemize}[nosep, leftmargin=1.4em]
  \item Did you snapshot \code{q.size()} \emph{before} the inner loop? (Trap 1.)
  \item Do both child pushes have non-null guards, and is the empty-tree case returned
        early? (Traps 2--3.)
  \item For a per-row aggregate, is your fold identity correct for negative values?
  \item Is your cosmetic change (zigzag/order) in the \emph{output}, not the queue?
        (\probref{c03-zigzag}'s dead end.)
\end{itemize}
\end{stuckbox}

\begin{practiceset}
  \practice{Binary Tree Level Order Traversal II}{LeetCode 107}
    {https://leetcode.com/problems/binary-tree-level-order-traversal-ii/}{LC Medium L2}
    {Run ordinary level order, then \code{reverse(res)} -- bottom-up is one line away from
    \coderef{c03-bfs}. (The plain top-down LC 102 is this template verbatim.)}
  \practice{Average of Levels in Binary Tree}{LeetCode 637}
    {https://leetcode.com/problems/average-of-levels-in-binary-tree/}{LC Easy L2}
    {Per-row fold of \code{sum}, then divide by the frozen \code{sz}; use \code{long long}
    or \code{double} to avoid integer truncation and overflow.}
  \practice{Populating Next Right Pointers in Each Node II}{LeetCode 117}
    {https://leetcode.com/problems/populating-next-right-pointers-in-each-node-ii/}{LC Medium L4}
    {Within each level, link \code{prev->next = cur} as you pop; the harder follow-up is to
    do it in $\bigO{1}$ space by treating the current level as a linked list that builds
    the next -- the ``BFS without a queue'' trick.}
\end{practiceset}

\begin{recallbox}[10-second refresher -- 3.3 Breadth-first / level order]
Queue, not stack: BFS sweeps one depth at a time. \emph{Freeze} \code{q.size()} to mark the
row, pop that many, enqueue children (guarded, non-null), and the frozen width is your
per-row loop bound. Per-row answers = same loop, different reducer; cosmetic reordering
lives in the output, never the frontier. Memory is $\bigO{w}$ -- no $\bigO{1}$ space here.
\selftest{
  \item Write \coderef{c03-bfs} from memory, marking the snapshot line and why it must be
        frozen.
  \item State BFS time and space, and the case where the queue hits $\bigO{n}$.
  \item On $R$, give level order, zigzag, and the per-row maximum without looking.
  \item Why must you not reorder child \emph{pushes} to get zigzag?
  \item Give the correct fold-identity for per-row max and for per-row sum.
}
\end{recallbox}
